% Part: first-order-logic
% Chapter: completeness
% Section: outline

\documentclass[../../../include/open-logic-section]{subfiles}

\begin{document}

\iftag{FOL}
      {\olfileid{fol}{com}{out}}
      {\olfileid{pl}{com}{out}}

\olsection{सिद्धतेचा आराखडा}

संपूर्णता प्रमेयाची सिद्धता काहीशी गुंतागुंतीची आहे आणि ती प्रथम वाचताना
वाट चुकणे सोपे आहे. म्हणून सिद्धतेचा आराखडा पाहू. पहिली पायरी म्हणजे
दृष्टिकोनातील बदल; त्यामुळे सिद्धतेकडे जाणारा मार्ग दिसू लागतो. संपूर्णतेचा
अर्थ “जेव्हा जेव्हा $\Gamma \Entails !A$, तेव्हा $\Gamma \Proves !A$”
असा घेतल्यास कल्पनाही सुचणे कठीण जाऊ शकते: कारण $\Gamma \Proves !A$ हे
दाखवण्यासाठी आपल्याला !!a{derivation} शोधावी लागते, आणि $\Gamma \Entails
!A$ हे गृहीतक त्यासाठी कोणत्याही प्रकारे मदत करते असे दिसत नाही. काही
सिद्धता-पद्धतींसाठी !!a{derivation} थेट रचता येते; पण आपण थोडा वेगळा मार्ग
घेऊ. आवश्यक दृष्टिकोन-बदल असा आहे: संपूर्णता पुढीलप्रमाणेही मांडता येते:
“$\Gamma$ सुसंगत असेल, तर तो पूर्ततायोग्य असतो.” कदाचित~$\Gamma$ मधील
माहिती आणि तो सुसंगत असल्याचे गृहीतक वापरून आपण~$\Gamma$ मधील प्रत्येक
\iftag{FOL}{!!{sentence}}{!!{formula}} पूर्ण करणारी
\iftag{FOL}{!!a{structure}}{!!a{valuation}} रचू शकू. शेवटी, आपल्याला
कोणत्या प्रकारची \iftag{FOL}{!!{structure}}{!!{valuation}} हवी आहे हे
माहीत आहे:~$\Gamma$ जिचे वर्णन करतो अशीच!{}

$\Gamma$ मध्ये फक्त \iftag{FOL}{आण्विक
  !!{sentence}s}{!!{propositional variable}s} असतील, तर त्याच्यासाठी
प्रतिमान रचणे सोपे आहे.\iftag{FOL}{ सर्व आण्विक !!{sentence}s ही
  $\Atom{P}{a_1,\dots,a_n}$ या रूपाची आहेत असे समजा; येथे $a_i$ ही
  !!{constant}s आहेत.}{} आपल्याला फक्त
\iftag{FOL}{%
  !!a{domain}~$\Domain{M}$ आणि~$P$ साठी असे निर्देशन ठरवायचे आहे की
  $\Sat{M}{\Atom{P}{a_1,\ldots,a_n}}$. पण ते फार कठीण नाही:
  $\Domain{M} = \Nat$, $\Assign{\Obj c_i}{M} = i$ असे ठेवा आणि प्रत्येक
  $\Atom{P}{a_1,\ldots,a_n} \in \Gamma$ साठी $\tuple{k_1,
    \dots, k_n}$ हे क्रमित बहुपद $\Assign{P}{M}$ मध्ये घाला; येथे $k_i$
  हा स्थिर चिन्ह~$a_i$ चा निर्देशांक आहे (म्हणजे $a_i \ident \Obj
  c_{k_i}$).}{%
  प्रत्येक $p \in \Gamma$ साठी $\pSat{v}{p}$ अशी !!a{valuation}~$\pAssign{v}$.
  त्यासाठी $\pAssign{v}(p) = \True$ तेव्हा आणि केवळ तेव्हाच ठेवू, जेव्हा
  $p \in \Gamma$.}

आता $\Gamma$ मध्ये~$\lnot !B$ हे एखादे !!{formula} असून $!B$ आण्विक आहे
असे समजा. $\iftag{FOL}{\Struct{M}}{\pAssign{v}}$ ची रचना $\lnot !B$ सत्य
करण्याच्या शक्यतेत अडथळा आणेल, अशी काळजी वाटू शकते. पण येथेच~$\Gamma$ ची
सुसंगतता उपयोगी पडते: $\lnot !B \in \Gamma$ असेल, तर $!B \notin \Gamma$;
अन्यथा $\Gamma$ विसंगत असेल. आणि $!B \notin \Gamma$ असेल, तर आपल्या
$\iftag{FOL}{\Struct{M}}{\pAssign{v}}$ च्या रचनेनुसार
$\iftag{FOL}{\Sat/{M}}{\pSat/{v}}{!B}$; म्हणून
$\iftag{FOL}{\Sat{M}}{\pSat{v}}{\lnot !B}$. आतापर्यंत सर्व ठीक आहे.

$\Gamma$ मध्ये गुंतागुंतीची, अनाण्विक सूत्रे असतील तर काय? उदाहरणार्थ,
त्यात $!A \land !B$ आहे असे समजा. ते सत्य करण्यासाठी $!A$ आणि $!B$ ही
दोन्ही~$\Gamma$ मध्ये असल्याप्रमाणे आपण पुढे गेले पाहिजे. आणि $!A \lor
!B \in \Gamma$ असेल, तर त्यांपैकी किमान एक सत्य करावे लागेल; म्हणजे
त्यांपैकी एक~$\Gamma$ मध्ये असल्याप्रमाणे पुढे जावे लागेल.

यावरून पुढील कल्पना सुचते: आपण~$\Gamma$ मध्ये अतिरिक्त !!{formula}s
अशा रीतीने भर घालतो की (a)~मिळणारा संच सुसंगत राहील, (b)~प्रत्येक
संभाव्य आण्विक !!{sentence}~$!A$ साठी मिळणाऱ्या संचात $!A$ किंवा
$\lnot !A$ यांपैकी एक असेल, आणि (c)~संचात $!A \land !B$ असेल तेव्हा
$!A$ व $!B$ ही दोन्ही असतील, तर $!A \lor !B$ असेल तेव्हा $!A$ किंवा
$!B$ यांपैकी किमान एकही असेल, इत्यादी. हे आपण करत राहतो (कदाचित
अनंतकाळपर्यंत). अशा रीतीने भर घातलेल्या सर्व !!{formula}s च्या संचाला
$\Gamma^*$ म्हणू. मग वरील रचनेतून आपल्याला अशी
\iftag{FOL}{!!a{structure}~$\Struct{M}$}{!!a{valuation}~$\pAssign{v}$}
मिळेल की ती~$\Gamma^*$ मधील सर्व वाक्ये पूर्ण करते, हे विगमनाने सिद्ध
करता येईल. $\Gamma \subseteq \Gamma^*$ असल्यामुळे ती~$\Gamma$ मधीलही
सर्व वाक्ये पूर्ण करेल. प्रत्यक्षात (a) आणि~(b) यांची हमी देणे पुरेसे
ठरते. ज्या वाक्यसंचासाठी (b) खरे असते त्याला \emph{संपूर्ण} म्हणतात.
त्यामुळे सुसंगत संच~$\Gamma$ चा विस्तार करून सुसंगत आणि संपूर्ण संच
$\Gamma^*$ मिळवणे हे आपले उद्दिष्ट असेल.

\iftag{FOL}{%
या योजनेत एक अडचण आहे: $\lexists[x][!A(x)] \in \Gamma$ असेल, तर या
प्रक्रियेत एखादे !!{constant}~$c$ निवडून $!A(c)$ ची भर घालता यावी अशी
आपली अपेक्षा आहे. पण ते नेहमी करता येईल हे कसे कळणार? कदाचित आपल्या
भाषेत मोजकीच !!{constant}s असतील आणि त्यांपैकी प्रत्येकासाठी
$\lnot !A(c) \in \Gamma$ असेल. आपण $!A(c)$ चीही भर घालू शकत नाही,
कारण त्यामुळे संच विसंगत होईल; आणि $\Struct{M}$ ने $!A(c)$ सत्य करावे की
$\lnot !A(c)$ हे आपल्याला कळणार नाही. शिवाय, $\Gamma$ मध्ये अशा भाषेतील
वाक्येच असू शकतात जिच्यात एकही स्थिर चिन्ह नाही (उदा., संचसिद्धान्ताची
भाषा).

या समस्येचा उपाय म्हणजे सुरुवातीलाच अनंत स्थिरे आणि त्यांना संख्यापकांशी
योग्य प्रकारे जोडणारी वाक्ये यांची भर घालणे. (अर्थात, त्यामुळे विसंगती
निर्माण होऊ शकत नाही, हे आपल्याला तपासावे लागेल.)

आण्विक वाक्यांमध्ये फक्त !!{constant}s असतील, तर आपली मूळ रचना चांगली
चालते. पण भाषेत !!{function}s सुद्धा असू शकतात. अशा वेळी सर्व काही योग्य
रीतीने चालण्यासाठी या !!{function}s ना नेमून द्यायची~$\Nat$ वरील योग्य
फलने शोधणे अवघड ठरू शकते. म्हणून आणखी एक युक्ती वापरू: $\Obj c_i$ चा
अर्थ लावण्यासाठी $i$ वापरण्याऐवजी !!{constant}s चा संचच प्रांत म्हणून
घेऊ. मग $\Struct M$ प्रत्येक !!{constant} ला तेच स्थिर नेमून देऊ शकते:
$\Assign{\Obj c_i}{M} = \Obj c_i$. पण पूर्ण मार्गच का घेऊ नये:
$\Domain{M}$ हा भाषेतील सर्व \emph{पदांचा} संच असू द्या!{} असे केल्यास
!!{function}s ना फलनांचे स्पष्ट निर्देशन मिळते (या फलनांची प्रचले पदे असतात
आणि त्यांची मूल्येही पदे असतात): !!{function}~$\Obj f^n_i$ ला आपण असे
फलन नेमून देतो की $n$ पदे $t_1$, \dots,~$t_n$ आदान म्हणून दिल्यावर ते
$\Obj f^n_i(t_1, \dots, t_n)$ हे पद मूल्य म्हणून देते.

कोड्याचा शेवटचा भाग म्हणजे~$\eq$ चे काय करायचे. !!{predicate}~$\eq$ चा
अर्थ स्थिर आहे: $\Sat{M}{\eq[t][t']}$ तेव्हा आणि केवळ तेव्हाच, जेव्हा
$\Value{t}{M} = \Value{t'}{M}$. आता पद~$t$ चे !!{value} हे $t$ स्वतःच
असेल अशी मांडणी केली, तर $t$ आणि $t'$ हे अगदी एकच पद असल्याखेरीज ही
!!{structure} $\eq[t][t']$ या रूपातील \emph{एकही} वाक्य सत्य करणार नाही.
आणि ही अर्थातच समस्या आहे; कारण !!{function}s असलेल्या भाषेतील जवळजवळ
प्रत्येक रोचक उपपत्तीमध्ये $t$ व $t'$ ही वेगळी पदे असूनही $\eq[t][t']$
या रूपाची वाक्ये प्रमेये असतील (उदा., अंकगणिताच्या उपपत्तींमध्ये
$\eq[(\Obj 0+ \Obj 0)][\Obj 0]$). ही समस्या सोडवण्यासाठी आपण
$\Struct M$ चा प्रांत बदलतो: $\Domain{M}$ मधील वस्तू म्हणून पदे वापरण्याऐवजी
पदांचे संच वापरतो; आणि प्रत्येक संचात~$\Gamma$ मधील वाक्यांनी समान असणे
आवश्यक ठरवलेली सर्व पदे असतात. उदाहरणार्थ, $\Gamma$ ही अंकगणिताची उपपत्ती
असेल, तर अशा संचांपैकी एकात $\Obj 0$, $(\Obj 0 + \Obj 0)$, $(\Obj
0 \times \Obj 0)$, इत्यादी असतील. हाच संच आपण $\Obj 0$ ला नेमून देऊ;
आणि तो त्यातील सर्व पदांचे मूल्यही आहे, उदा., $(\Obj 0 + \Obj 0)$ चेही,
हे निष्पन्न होईल. म्हणून सुधारित !!{structure} मध्ये
$\eq[(\Obj 0+ \Obj 0)][\Obj 0]$ हे वाक्य सत्य असेल.}{}

मग आपण पुढीलप्रमाणे काम करू. प्रथम !!{complete} सुसंगत संचांच्या गुणधर्मांचा
अभ्यास करू. विशेषतः, !!a{complete} सुसंगत संचात $!A \land !B$ तेव्हा आणि
केवळ तेव्हाच असते, जेव्हा त्यात $!A$ व~$!B$ ही दोन्ही असतात; $!A \lor !B$
तेव्हा आणि केवळ तेव्हाच असते, जेव्हा त्यांपैकी किमान एक असते; इत्यादी,
हे सिद्ध करू (\olref[ccs]{prop:ccs}).\iftag{FOL}{ त्यानंतर वाक्यांच्या
  “संतृप्त” संचांची व्याख्या करून त्यांचा अभ्यास करू. प्रत्येक संख्यकीकृत
  !!{sentence} ला तिच्या दृष्टांतांशी जोडणाऱ्या सशर्तांचा समावेश असलेला
  संच म्हणजे संतृप्त संच (\olref[hen]{defn:henkin-exp}). कोणत्याही सुसंगत
  संच~$\Gamma$ चा विस्तार करून संतृप्त संच~$\Gamma'$ नेहमी मिळवता येतो,
  हे आपण दाखवू (\olref[hen]{lem:henkin}). एखादा संच सुसंगत, संतृप्त आणि
  !!{complete} असेल, तर त्यात \iftag{prvEx}{$\lexists[x][!A(x)]$ तेव्हा
    आणि केवळ तेव्हाच असते, जेव्हा एखाद्या बंद पद~$t$ साठी त्यात $!A(t)$
    असते}{}\iftag{defEx,defAll}{}{ आणि
  }\iftag{prvAll}{$\lforall[x][!A(x)]$ तेव्हा आणि केवळ तेव्हाच असते,
    जेव्हा प्रत्येक बंद पद~$t$ साठी त्यात $!A(t)$ असते}{}
  (\olref[hen]{prop:saturated-instances}), हाही गुणधर्म त्याला असतो.}{}
त्यानंतर \iftag{FOL}{संतृप्त}{} सुसंगत संच
$\iftag{FOL}{\Gamma'}{\Gamma}$ घेऊन त्याचा विस्तार करून
\iftag{FOL}{संतृप्त, सुसंगत आणि}{सुसंगत आणि} !!{complete} संच~$\Gamma^*$
मिळवता येतो, हे दाखवू (\olref[lin]{lem:lindenbaum}). याच~$\Gamma^*$ चा
वापर आपण आपले \iftag{FOL}{पद-प्रतिमान~$\Struct M(\Gamma^*)$}{!!{valuation}~
$\pAssign v(\Gamma^*)$} परिभाषित करण्यासाठी करू.
\iftag{FOL}{पद-प्रतिमानाचा प्रांत हा बंद पदांचा संच असतो आणि त्यातील
  !!{predicate}s चा अर्थ आण्विक !!{sentence}s}{मूल्यन हे !!{propositional
    variable}s} यांच्या~$\Gamma^*$ मधील सदस्यत्वाने ठरते
(\olref[mod]{defn:termmodel}). \iftag{FOL}{संतृप्त,}{} संपूर्ण सुसंगत
संचांचे गुणधर्म वापरून
$\iftag{FOL}{\Sat{M(\Gamma^*)}}{\pSat{v(\Gamma^*)}}{!A}$ तेव्हा आणि
केवळ तेव्हाच, जेव्हा $!A \in \Gamma^*$, हे दाखवू
(\olref[mod]{lem:truth}); म्हणून विशेषतः
$\iftag{FOL}{\Sat{M(\Gamma^*)}}{\pSat{v(\Gamma^*)}}{\Gamma}$.
\iftag{FOL}{शेवटी, $\Gamma$ मध्ये~$\eq$ सुद्धा असेल तर पद-प्रतिमानाची
  व्याख्या कशी करायची याचा विचार करू
  (\olref[ide]{defn:term-model-factor}) आणि ते~$\Gamma^*$ ची पूर्ति करते,
  हे दाखवू (\olref[ide]{lem:truth}).}{}

\end{document}
